\section{Centered 65/70/73 Degree STL Box-Receiver Result}

The new STL file
\texttt{panels\_only\_rectangular\_3floors\_250200250\_657073deg.stl} was first tested
using the same transform as the older \texttt{606570} STL.  That first result was too low
and strongly asymmetric, indicating that the new mesh had a different CAD origin.  After
checking the mesh bounds, the reflector was recentered over the box receiver using the
translation
\[
    \mathbf{t} =
    \begin{bmatrix}
        -0.2005135 \\
        -0.3374645 \\
        -0.2801248
    \end{bmatrix}
    \mathrm{m}.
\]
The corrected scene was then rerun with \SI{10000000} rays, DNI
\SI{1000}{W.m^{-2}}, mirror reflectance 0.85, slope error \SI{2}{mrad}, and the fixed
\SI{0.30}{m} wide box receiver.

\begin{table}[h]
\centering
\caption{Centered 10M-ray result for the 250/200/250 mm, 65/70/73 degree STL.}
\begin{tabular}{lrr}
\toprule
Face & Power (W) & Peak flux (\si{W.m^{-2}}) \\
\midrule
Glass top & 361.113 & 8977.013 \\
Bottom / battery proxy & 324.676 & 5205.449 \\
North wall & 9.099 & 1283.669 \\
South wall & 9.122 & 1250.704 \\
East wall & 9.101 & 1263.920 \\
West wall & 9.115 & 1273.035 \\
\bottomrule
\end{tabular}
\end{table}

The centered run gives a total absorbed power of
\[
    P_{\mathrm{abs}} = \SI{361.113}{W},
\]
with
\[
    P_{\mathrm{bottom}} = \SI{324.676}{W}
\]
absorbed on the bottom face, which is used here as the battery proxy.  This means
approximately
\[
    \frac{324.676}{361.113}\times 100\% \approx 89.91\%
\]
of the absorbed power reaches the intended battery-facing surface.  The nearly equal wall
powers, all about \SI{9.1}{W}, confirm that the corrected STL placement is well centered.

\begin{figure}[h]
    \centering
    % Replace this placeholder with the rendered flux map when compiling the final report.
    \includegraphics[width=0.95\linewidth]{images/unfolded_box_flux.png}
    \caption{Placeholder for the unfolded flux map of the centered 10M-ray run.}
    \label{fig:centered-657073-flux}
\end{figure}
